\section{Hysteres et fluage}


\subsection{Description du dispositif}

Le dispositif expérimental (Figure 2 du protocole) comprend :
\begin{itemize}
    \item Un \textbf{laser He-Ne} ($\lambda = 632.8$ nm) projeté sur une \textbf{lame flexible miroir} (Invar).
    \item Un \textbf{actuateur piézo-électrique} (P-117-00, déplacement max = 82 $\mu$m pour 150 V).
    \item Une \textbf{photodiode 4 quadrants} pour mesurer la déviation du faisceau.
    \item Un \textbf{générateur de fonctions} (signaux triangulaires/carrés) et un \textbf{oscilloscope USB PicoScope}.
\end{itemize}

Le déplacement du piézo courbe la lame, déviant le laser. La photodiode convertit cette déviation en signal électrique, permettant de tracer $d(V)$.

\begin{figure}[H]
    \centering
    \includegraphics[scale=0.4]{schema_hysterese.png}
    \label{fig:schema_hysterese}
    \caption{Schema du dispositif de mesure de l'hysterese}
\end{figure}



\subsection{Théorie : Hystérèse et fluage}
\subsubsection{Hystérèse}
L'hystérèse dans les matériaux piézo-électriques résulte de :

\begin{itemize}
    \item La \textbf{réorientation des domaines ferroélectriques} sous champ électrique.
    \item Les \textbf{frottements internes} et inerties mécaniques.
\end{itemize}
Le coefficient de non-linéarité est calculé comme :

\begin{equation}
    \text{Non-linéarité (\%)} = \frac{\text{Écart max (C-D)}}{\text{Plage totale (A-B)}} \times 100
    \label{eq:eq_lin}
\end{equation}

où A-B = déplacement total, C-D = écart maximal par rapport à la linéarité (Figure 1).

\subsubsection{Fluage (Creep)}
Le fluage décrit la \textbf{réponse temporelle} du piézo à un saut de tension. Il est modélisé par :

\begin{equation}
    y(t) = a \left(1 - e^{-t/\tau}\right) + b \quad \text{(montée)}
    \label{eq:t_m}
\end{equation}

\begin{equation}
    y(t) = a e^{-t/\tau} + b \quad \text{(descente)}
    \label{eq:t_d}
\end{equation}

où $\tau$ est la constante de temps (typiquement de l'ordre de la ms à la s).

\subsection{Resultats et Analyse}

\subsubsection{Hysterese}

Chaque courbe d'hysterese (voir Fig. \ref{fig:1V_hyst}) nous permet de tracer d(V) (voir Fig. \ref{fig:1V_lin}). Pour 1V sur les graphiques susmentionnés, nous constatons que la pente diminue en fonction de la fréquence (ou l'inverse).


C'est également le cas pour 2V comme nous l'observer sur Fig. \ref{fig:2V_lin} et pour 3V sur Fig. \ref{fig:3V_lin}.

\begin{figure}[H]
    \centering
    \includegraphics[scale=0.5]{1_1V.png}
    \label{fig:2V_lin}
    \caption{d(V) pour 2V}
\end{figure}

\begin{figure}[H]
    \centering
    \includegraphics[scale=0.5]{2_2V.png}
    \label{fig:2V_lin}
    \caption{d(V) pour 2V}
\end{figure}

\begin{figure}[H]
    \centering
    \includegraphics[scale=0.5]{3_3V.png}
    \label{fig:3V_lin}
    \caption{d(V) pour 3V}
\end{figure}


Or, fort est de constater que plus la tension est élevée, plus la différence entre les pentes est élevée entre les différentes fréquences. Nous pouvons observer les différentes pentes sur Fig. \ref{fig:pentes} d'après les valeurs dans Tab. \ref{tab:pentes} dont le pourcentage de non-linearite est donné par ...
\begin{table}[h]
\centering
\caption{Résultats des mesures d'hystérésis pour différentes tensions et fréquences}
\label{tab:hysteresis_results}
\begin{tabular}{|c|c|c|c|c|c|c|}
\hline
\textbf{Fichier} & \textbf{A} & \textbf{B} & \textbf{C} & \textbf{D} & \textbf{AB (V)} & \textbf{CD (V)} & \textbf{Non-linéarité} \\
\hline
1V5Hz\_15.csv & 0.013 & 0.036 & (1.453, 0.0274) & (1.453, 0.0201) & 0.022 & 0.007 & 0.328 \\
\hline
1V10Hz\_15.csv & 0.016 & 0.036 & (1.505, 0.0316) & (1.505, 0.0210) & 0.020 & 0.011 & 0.519 \\
\hline
1V20Hz\_15.csv & 0.017 & 0.032 & (1.526, 0.0310) & (1.526, 0.0194) & 0.015 & 0.012 & 0.760 \\
\hline
2V5Hz\_15.csv & 0.003 & 0.049 & (1.563, 0.0353) & (1.563, 0.0179) & 0.046 & 0.017 & 0.381 \\
\hline
2V10Hz\_15.csv & 0.005 & 0.041 & (1.679, 0.0364) & (1.679, 0.0157) & 0.036 & 0.021 & 0.571 \\
\hline
2V20Hz\_15.csv & 0.011 & 0.039 & (1.549, 0.0363) & (1.549, 0.0143) & 0.028 & 0.022 & 0.783 \\
\hline
3V5Hz\_15.csv & -0.003 & 0.053 & (1.735, 0.0405) & (1.735, 0.0167) & 0.056 & 0.024 & 0.421 \\
\hline
3V10Hz\_15.csv & -0.001 & 0.049 & (1.833, 0.0427) & (1.833, 0.0133) & 0.050 & 0.029 & 0.587 \\
\hline
3V20Hz\_15.csv & 0.004 & 0.043 & (1.614, 0.0389) & (1.614, 0.0094) & 0.039 & 0.030 & 0.765 \\
\hline
\end{tabular}
\end{table}

\begin{figure}[H]
    \centering
    \includegraphics[scale=0.5]{.png}
    \label{fig:3V_lin}
    \caption{d(V) pour 3V}
\end{figure}



\subsubsection{Fluage}
\

En un second temps, nous imposons un signal carré au système. Nous mesurons la réponse électrique de l'échantillon piezo comme nous pouvons l'observer sur Fig.\ref{fig:fluage}. 
\begin{figure}[H]
    \centering
    \includegraphics[scale=0.5]{fluage_fit_montée_descente.png}
    \label{fig:fluage}
    \caption{Courbe de fluage}
\end{figure}

Nous determinons les temps de montée et de descente sur Fig. ... . Nous constatons que le temps de montée est plus bas que le temps de descente (tau1 > tau2). 

\begin{table}[h]
\centering
\caption{Résultats de l'ajustement exponentiel pour la montée et la descente}
\begin{tabular}{|c|c|c|}
\hline
\textbf{Paramètre} & \textbf{Montée} & \textbf{Descente} \\
\hline
$A$ & $8.651228 \times 10^{-2} \pm 4.91 \times 10^{-5}$ & $8.352762 \times 10^{-2} \pm 3.15 \times 10^{-5}$ \\
\hline
$\tau$ (s) & $1.051606 \times 10^{-2} \pm 1.44 \times 10^{-5}$ & $1.078764 \times 10^{-2} \pm 1.00 \times 10^{-5}$ \\
\hline
$B$ & $-3.102036 \times 10^{-2} \pm 5.33 \times 10^{-5}$ & $-1.933693 \times 10^{-2} \pm 1.82 \times 10^{-5}$ \\
\hline
\end{tabular}
\label{tab:resultats_fluage}
\end{table}

Cela s'explique par le fait que ...


